\subsubsection{Tangens als Quotient}

Aus den Definitionen der trigonometrischen Funktionen folgt unmittelbar:

\begin{KeyPointBox}{Tangens als Quotient}{fundamentals_tangent_quotient}
  \noindent
  \begin{minipage}[c]{0.73\linewidth}
    Im Quotienten aus Sinus und Kosinus kürzt sich die Hypotenuse $c$ heraus:
    \begin{ClozeArea}[23]{6}
      \[
        \frac{\sin(\alpha)}{\cos(\alpha)}
          =\frac{a/c}{b/c}
           =\frac{a}{b}
           =\tan(\alpha)
        \quad\Longrightarrow\quad
        \tan(\alpha)=\frac{\sin(\alpha)}{\cos(\alpha)}.
      \]
    \end{ClozeArea}
  \end{minipage}\hfill
  \begin{minipage}[c]{0.23\linewidth}
    \centering
    \CachedTikzFigure{content/resources/chapters/fundamentals/tikz/right_triangle_trigonometric_ratios_example.tex}
  \end{minipage}
\end{KeyPointBox}

Damit können wir nun die Werte der Tangensfunktion $\tan(\ldots)$ aus den entsprechenden Sinus- und
Kosinuswerten berechnen.

\IfClozeTF{\clearpage}{}
